\documentclass[11pt]{article}

\usepackage[margin=1in]{geometry}
\usepackage{amsmath,amssymb,amsthm}
\usepackage{graphicx}
\usepackage[colorlinks=true,linkcolor=blue,citecolor=blue,urlcolor=blue]{hyperref}

\newtheorem{theorem}{Theorem}
\newtheorem{remark}{Remark}
\newcommand{\N}{\mathbb{N}}
\newcommand{\C}{\mathbb{C}}
\newcommand{\linf}{\ell_\infty}
\newcommand{\cinf}{c_0}
\newcommand{\chiA}{\chi_A}

\title{Papathanasiou's dense generators are not C*-minimal}
\author{Agent packet for run fa\_banach\_001}
\date{2026-07-01}

\begin{document}
\maketitle

\section*{Status}

This is a claimed full negative answer, subject to human review, to the
minimality question in Section 4 of Franca--Garces \cite{FrancaGarces}.

\section*{Source passage}

The source paper constructs a copy \(A_3\cong \linf\) inside
\((\linf\setminus \cinf)\cup\{0\}\).  In the final considerations, after
constructing another subalgebra \(A_4\) with a minimal generator system of
cardinality \(\mathfrak c\), it says:

\begin{quote}
Although it is even possible to find a dense subalgebra of \(A_3\) as in
\cite{Papathanasiou}, we do not know if that system of generators is minimal
when generating \(\ell_{\infty}\) (or \(A_3\)) as a C*-algebra.
\end{quote}

The relevant source crop is included as Figure~\ref{fig:source-crop}.

\begin{figure}[h]
\centering
\includegraphics[width=\textwidth]{figures/open_problem_crop.png}
\caption{The source passage on page 17 of the rendered arXiv source.}
\label{fig:source-crop}
\end{figure}

\section*{Papathanasiou's family}

In \cite[proof of Theorem 2]{Papathanasiou}, for every
\(A\subseteq \N\) and \(j\in\N\) there is a bounded sequence
\(f_{(A,j)}\) with \(\|f_{(A,j)}\|\leq 1/j\), and
\[
        h_{(A,j)}=\chi_A+f_{(A,j)}.
\]
Hence \(h_{(A,j)}\to \chi_A\) in the sup norm as \(j\to\infty\), for each
fixed \(A\).  The generated algebra is dense in \(\linf\), and the proof in
\cite{Papathanasiou} shows that it lies in
\((\linf\setminus \cinf)\cup\{0\}\) and is generated by an algebraically
independent family.  The same argument is available over the complex scalars,
as noted in \cite[Remark 3]{Papathanasiou}.

\section*{Proof intuition}

The density mechanism has one redundant coordinate built into it.  For each
set \(A\subseteq \N\), the construction does not use one approximant
\(h_{(A,j)}\), but an entire tail converging to \(\chi_A\).  Removing a
single generator removes at most one term from one of these tails.  All
characteristic functions \(\chi_A\) therefore remain in the C*-closure of the
surviving family, and those characteristic functions already generate
\(\ell_\infty\) as a C*-algebra.

\begin{theorem}
Let
\[
        H=\{h_{(A,j)}: A\subseteq \N,\ j\in\N\}\subseteq \linf
\]
be Papathanasiou's generator family.  For every
\((A_0,j_0)\in \mathcal P(\N)\times\N\),
\[
        C^*\bigl(H\setminus\{h_{(A_0,j_0)}\}\bigr)=\linf .
\]
Consequently \(H\) is not a minimal C*-generating set for \(\linf\).  If
\(\Phi:\linf\to A_3\) is any *-isomorphism onto the copy \(A_3\) appearing in
\cite{FrancaGarces}, then \(\Phi(H)\) is not a minimal C*-generating set for
\(A_3\).
\end{theorem}

\begin{proof}
Fix \((A_0,j_0)\), set
\[
        H_0=H\setminus\{h_{(A_0,j_0)}\},
\]
and let \(B=C^*(H_0)\), the C*-subalgebra of \(\linf\) generated by \(H_0\).

We first show that \(B\) contains every characteristic function \(\chi_A\),
\(A\subseteq\N\).  If \(A\neq A_0\), then \(h_{(A,j)}\in H_0\) for all
\(j\), and since \(\|h_{(A,j)}-\chi_A\|\leq 1/j\), closedness of \(B\) gives
\(\chi_A\in B\).  If \(A=A_0\), choose any sequence \(j_n\to\infty\) with
\(j_n\neq j_0\) for all \(n\).  Then \(h_{(A_0,j_n)}\in H_0\) and
\[
        \|h_{(A_0,j_n)}-\chi_{A_0}\|\leq 1/j_n\to 0,
\]
so again \(\chi_{A_0}\in B\).

Thus \(B\) contains \(\chi_A\) for every \(A\subseteq\N\).  Every finite-valued
bounded sequence on \(\N\) is a finite linear combination of such
characteristic functions, so all finite-valued bounded sequences lie in
\(B\).  Finally, finite-valued bounded sequences are uniformly dense in
\(\ell_\infty\): given \(x\in\ell_\infty\) and \(\varepsilon>0\), partition
the bounded range of \(x\) into finitely many subsets of diameter less than
\(\varepsilon\), and approximate \(x\) by the corresponding step sequence.
Since \(B\) is norm closed, \(B=\ell_\infty\).

This proves that deleting the single generator \(h_{(A_0,j_0)}\) still leaves
a C*-generating family for \(\ell_\infty\).  As \((A_0,j_0)\) was arbitrary,
no element of \(H\) is indispensable, so \(H\) is not minimal.

If \(\Phi:\ell_\infty\to A_3\) is a *-isomorphism, then
\[
 C^*\bigl(\Phi(H)\setminus\{\Phi(h_{(A_0,j_0)})\}\bigr)
 =\Phi\!\left(C^*(H\setminus\{h_{(A_0,j_0)}\})\right)
 =\Phi(\ell_\infty)=A_3.
\]
Hence the transferred system in \(A_3\) is not minimal either.
\end{proof}

\section*{Verification notes}

The proof uses only the displayed approximation
\(\|h_{(A,j)}-\chi_A\|\leq 1/j\), the fact that the index set contains all
pairs \((A,j)\), and the elementary density of finite-valued sequences in
\(\ell_\infty\).  No set-theoretic assumptions are used.  The argument is
unchanged after passing through a *-isomorphism \(\ell_\infty\to A_3\).

\section*{Novelty and search bounds}

Before promotion, the lightweight run indexes were searched for arXiv id
2508.19174, the source title, ``compact multipliers'', ``algebrability'',
``Papathanasiou'', ``dense subalgebra'', and generator-minimality phrases.
No existing packet or attempt answered this target.  A bounded web search on
2026-07-01 for the exact source sentence, the source title, and close phrases
combining Papathanasiou, \(\ell_\infty\setminus c_0\), minimal generators,
and C*-algebras did not surface an exact duplicate in returned snippets.

\section*{Human review recommendation}

The only point requiring review is whether the phrase ``that system of
generators'' in \cite{FrancaGarces} refers exactly to the family
\(\{h_{(A,j)}\}\) from \cite{Papathanasiou}, or to its image under a fixed
*-isomorphism \(\ell_\infty\to A_3\).  The theorem covers both readings.

\begin{thebibliography}{9}

\bibitem{FrancaGarces}
W. Franca and Jorge J. Garces,
\emph{Applications of compact multipliers to algebrability of
\((\ell_{\infty}\setminus c_0)\cup\{0\}\) and
\((B(\ell_2(\mathbb N))\setminus K(\ell_2(\mathbb N)))\cup\{0\}\)},
arXiv:2508.19174v1, 2025.

\bibitem{Papathanasiou}
Dimitris Papathanasiou,
\emph{Dense lineability and algebrability of
\(\ell^\infty\setminus c_0\)},
Proc. Amer. Math. Soc. 150 (2022), 991--996; arXiv:2102.03199.

\end{thebibliography}

\end{document}
